\section{DPM-Solver++ y mejoras prácticas}

\subsection{De la predicción de ruido a la predicción de datos}

La fórmula original de DPM-Solver se basa en la predicción de ruido $\boldsymbol{\epsilon}_\theta$. DPM-Solver++ propone una mejora clave: cambiar hacia la aproximación polinómica de la \textbf{predicción de datos} ($\hat{\mathbf{x}}_0$).

Dado que $\mathbf{x}_t = \alpha_t \mathbf{x}_0 + \sigma_t \boldsymbol{\epsilon}$, se tiene:
$$\hat{\mathbf{x}}_0(\mathbf{x}_t, t) = \frac{\mathbf{x}_t - \sigma_t \boldsymbol{\epsilon}_\theta(\mathbf{x}_t, t)}{\alpha_t}$$

\begin{importantbox}{Mejora central de DPM-Solver++}
DPM-Solver++ cambia el objeto de integración de $\boldsymbol{\epsilon}_\theta$ a $\hat{\mathbf{x}}_0$. Aunque son matemáticamente equivalentes, en un número finito de pasos, las variaciones de $\hat{\mathbf{x}}_0$ suelen ser más suaves que las de $\boldsymbol{\epsilon}_\theta$, lo que resulta en menores errores de aproximación polinómica. Los experimentos muestran que DPM-Solver++ es significativamente mejor que el DPM-Solver original con menos pasos (por ejemplo, entre 10 y 15 pasos).
\end{importantbox}

\begin{knowledgebox}{¿Por qué la predicción de datos es más suave?}
Intuitivamente, $\hat{\mathbf{x}}_0$ es una "estimación de la imagen limpia después de eliminar el ruido", que converge gradualmente a la imagen final durante el proceso de muestreo. En contraste, $\boldsymbol{\epsilon}_\theta$ (la estimación del ruido) cambia con mayor intensidad a lo largo de los pasos temporales, especialmente en las zonas donde el nivel de ruido varía rápidamente. Por lo tanto, realizar una aproximación polinómica de bajo orden para $\hat{\mathbf{x}}_0$ es más razonable.
\end{knowledgebox}

\subsection{Métodos multietapa (Multistep Method)}

Además de los métodos de un solo paso (donde cada punto intermedio se calcula desde cero), DPM-Solver++ también soporta \textbf{métodos multietapa}: utiliza las predicciones de ruido ya calculadas en pasos anteriores para construir aproximaciones de orden superior, sin necesidad de evaluaciones adicionales de la red.

Por ejemplo, el DPM-Solver++ multietapa de segundo orden:
\begin{enumerate}
    \item En el paso $i$, se tienen $\boldsymbol{\epsilon}_\theta(\mathbf{x}_{t_i}, t_i)$ y $\boldsymbol{\epsilon}_\theta(\mathbf{x}_{t_{i-1}}, t_{i-1})$.
    \item Se utiliza una extrapolación lineal de estos dos valores para aproximar la variación del ruido dentro del intervalo.
    \item Se sustituye en la fórmula de integral exponencial para calcular $\mathbf{x}_{t_{i+1}}$.
\end{enumerate}

\begin{importantbox}{Ventajas de los métodos multietapa}
Los métodos multietapa requieren solo 1 evaluación de red por paso (NFE=1), pero pueden alcanzar una precisión de orden 2 e incluso 3. El costo implica la necesidad de almacenar el búfer de predicciones de ruido de los pasos anteriores. Dado un presupuesto fijo de pasos, los métodos multietapa suelen ser la opción con mejor relación calidad-precio.
\end{importantbox}

\subsection{Estrategias de programación de pasos temporales}

El rendimiento de DPM-Solver también está influenciado por la estrategia de selección de los pasos temporales. Las estrategias comunes incluyen:

\begin{itemize}
    \item \textbf{Pasos temporales uniformes}: dividir equidistantemente en $t \in [0, T]$.
    \item \textbf{Pasos $\lambda$ uniformes}: dividir equidistantemente en $\lambda \in [\lambda_T, \lambda_0]$ — esto suele ser mejor porque la uniformidad en el espacio $\lambda$ se alinea mejor con la distribución del error numérico de las EDO.
    \item \textbf{Paso adaptable}: ajustar dinámicamente el paso según la estimación del error (como en los métodos DOPRI5).
\end{itemize}

\begin{warningbox}{Una mala selección de pasos temporales reducirá gravemente el rendimiento}
Al dividir uniformemente en el espacio $t$, el tamaño del paso en las regiones de bajo ruido (cuando $t$ se acerca a 0) resulta excesivamente grande en el espacio $\lambda$, lo que concentra el error de aproximación en los últimos pasos —y esto es exactamente la etapa crítica donde se determinan los detalles de la imagen. Por lo tanto, dividir uniformemente en el espacio $\lambda$ suele ser una elección por defecto más segura.
\end{warningbox}

\subsection{Resumen de la sección}

DPM-Solver++ mejora la precisión en un número finito de pasos al realizar aproximaciones polinómicas sobre la predicción de datos (en lugar de la predicción de ruido) e introduce métodos multietapa para lograr una precisión de orden superior con el costo de 1 NFE por paso. La programación de los pasos temporales (especialmente la división uniforme en el espacio $\lambda$) también tiene un impacto importante en el rendimiento final.
